\section{Combining local estimates in the row sum}
\label{sec:row-sums}

For a fixed pair of words, each reduced Gram label contributes a fibre of integral terms. We sum the squared-norm loss in \eqref{eq:row-l2} and, on the cells where it is available, the pointwise loss in \eqref{eq:row-abs}. We combine the
height bounds with the counts from Section~\ref{sec:flag-counts}
and separate the contributions according to their Gram
invariants. The exponents in the resulting global bounds are
used in the full Gram estimate to control the off-diagonal
part.

Fix two sequences \(u,u'\) and use their word norms \(T,T'\).
Write \(d=\dim\pi_u\), \(p=r+1\), \(m=2p\), \(h=cr^2\)
with the height constant \(c\ge1\), and \(f\) for the number of real
places. At a real place \(v\), write \(d_v\) for the dimension of the local factor of the definite representation \(\pi_u\). The common
archimedean scale and normalized amplitudes are
\[
              \upsilon\ge1,\qquad A_t=a_t/\sqrt d .
\]

Only terms with positive local Gram invariants \(G_{Z,v}\) enter this
sum. The pair of terms can still have a degenerate determinant in the
count of Section~\ref{sec:flag-counts}. A third distinction within
one term is whether its preceding minor \(\Delta_r\) vanishes. This
determines which local bound is available. Put \(g=g_0g_1\) and
\[
 G_Z=g^{-mf}\prod_vG_{Z,v},\quad N_T=T^f,\quad
 n_Z=\left(\frac{e\pi}{k\upsilon}\right)^{kf}(G_Z/N_T)^k .
\]
Choose \(C_h\) as in Lemma~\ref{lem:minor-orbital-iterate}
and write \(B=b^{C_hm^2}\). For leading minors \(D_{l,v}\) of the
reduced label define
\[
 s_v=\prod_{l=1}^{p}\max\{1,|D_{l,v}|/B\},\qquad
 \sigma_Z=\prod_vs_v,\qquad
 P=b^{(c^2+2\max(C_h-c,0))m^3 f}.
\]
The reduced height and the term volume satisfy
\(H_Z\le P\sigma_Z^4\) and \(\operatorname{vol}(t)\le H_Z^h\).

At a block end write \(d_{l,v}=|D_{l,v}|>0\), and put
\(G_v=G_{Z,v}\). The local squared-norm and pointwise losses, with
the fixed comparison constant \(C\), are respectively
\[
 \begin{aligned}
 L_v&=b^{Cm^4}\max(G_v,G_v^{-1})^{Cm^2}
           \prod_{l\ {\rm end}}\max(d_{l,v},d_{l,v}^{-1})^{Cm},\\
 P_v&=\left(b^{Cm}\max\big(\{1\}\cup
                 \{d_{l,v}^{-1}:l\text{ a block end}\}\big)\right)^{Cr^3}.
 \end{aligned}
\]
For a moment order \(q\) in the range of Lemma~\ref{lem:row-moment-budget}, set \(a=q/p\),
\(E=\exp(C_2m^6 f\log b)\), and \(\Gamma_S=b^{C_\Gamma m^5f}\).
Put \(t_v=T_{{\rm crit},v}/T'\), \(r_{DD}(u')=d_{u'}/D_*\) and \(\tau=(1+\delta)r_{DD}(u')\). The estimate on their product comes from the full marked dimension.

Write \(R=k+m-1\) for the ambient row rank. Multiplying the defining formulas for \(T_{{\rm crit},v}^k\) in Section~\ref{sec:local-amplitudes} yields exactly the archimedean, discriminant and level powers of \(D_*/d_{u',*}\) in the proof of Proposition~\ref{prop:prime-target-size}. The factors left over are
\[
 Z=\sum_{j=k+1}^R\log L_j+
        f\sum_{i=0}^1\sum_{j=k+1}^R\log(1-g_i^{-j}).
\]
In fact the archimedean factors at \(v\) are \(\kappa^{-[(m-1)k+m(m-1)/2]}\) times the ratio of the ambient and column Vandermondes. Proposition~\ref{prop:marked-column-dimension} now gives
\[
 \frac{D_*}{d_{u'}}=e^Z (1-(T')^{-k})^{-f}\prod_v t_v^k .
\]

For \(j\ge2\), \(\log\zeta(j)\le\sum_{n\ge2}n^{-j}\le3\cdot2^{-j}\), by integral comparison. The Euler bounds in Proposition~\ref{prop:prime-target-size} imply \(|Z|\le9f2^{-k}\). Also \(-\log(1-(T')^{-k})\le2(T')^{-k}\). Thus on the selected range
\[
 \prod_v t_v^{-k}\le e^{11f2^{-k}}\frac{d_{u'}}{D_*}\le\tau .
\]
The last inequality holds once \(11f2^{-k}\le\log(1+\delta)\), uniformly in the word. In particular the scale conditions \(f\le k\), \(k\ge m^5(\log b)x^2\) allow the tolerance \(\delta=b^{-7}\) used in the row estimates. \label{cor:row-sums-1}

\begin{lemma}[Products of local height losses]\label{lem:row-height-products}
For the reduced label and its common block ends, suppose
\(C_2\ge\max(6C,C,C^2+CC_h)\) and \(C+6CC_h\le C_\Gamma m\).
The local factors satisfy
\[
 \prod_v L_v\le\Gamma_S\sigma_Z^{C_2m^2},\qquad
 \prod_v P_v\le E\sigma_Z^{C_2m^4}.
\]
The second bound also applies to a product over a subset of places.
\end{lemma}

\begin{proof}
Put \(X=B^f\sigma_Z\). For any leading minor, multiplication of
its local bounds gives \(\prod_v\max(1,|D_{l,v}|)\le X\).
At a block end integrality gives \(\prod_v d_{l,v}\ge1\). Likewise
\(\prod_v G_v\ge1\) and \(G_v\le |D_{p,v}|^2\).
For a positive \(y\), both inverse-height factors are controlled by
\[
 \max(1,y^{-1})=\max(1,y)/y,\qquad
 \max(y,y^{-1})=\max(1,y)^2/y.
\]
There are at most \(m\) common ends. Thus the two products are at
most \(b^{Cm^4 f}X^{6Cm^2}\) and \((b^{Cmf}X^m)^{Cr^3}\),
respectively. The indicated inequalities on the constants absorb the
powers of \(B\). Each \(P_v\ge1\), which gives the subset assertion.
\end{proof}

When \(\Delta_r=0\) set \(K_Z=1\). Otherwise define the Cayley factor by
\[
 \mathrm{sh}_v^2=\frac{G_{Z,v}}{TT'|\Delta_{r,v}|^2},\qquad
 \mathrm{ch}_v=\sqrt{1+\mathrm{sh}_v^2},\qquad
 K_Z=\prod_v(1+\mathrm{ch}_v)/\mathrm{sh}_v.
\]
It satisfies \(K_Z^2\le9^f H_Z^3\): use
\(1+\sqrt{1+x^2}\le3\max(1,x)\) for \(x>0\) and the height bounds
on the norms, the preceding minor and the inverse Gram product.

\begin{auxlemma}[Evaluation of the cell norm at a positive Gram point]
\label{lem:unit-cell-evaluation}
Let \(J_\mu\) be the amplitude in Lemma~\ref{lem:minor-orbital-shift},
\(\mu=\lambda+\nu\), \(\|\nu\|_1\le C_nkm^3\).
Use the actual ordered words, compact types and root measures at a
marked real place. Write \(w(\gamma)\) for the orbital Jacobian of
the block containing \(\gamma\), and \(G=G_{Z,v}>0\).
For all sufficiently late marked rows, uniformly for every nonzero
pair of periods in this window,
\begin{equation}\label{eq:unweighted-unit-cell-evaluation}
 w(\gamma)\,\frac{\|J_\mu(\gamma)\|_{\rm HS}^2}{M_\mu}
       \le b^{C_bm^4}\max(1,G^{-1})^{C_bm^2},
 \qquad M_\mu=\|J_\mu\|_{\rm cell}^2 .
\end{equation}
The constant \(C_b\) depends on the fixed weight, character and
measure bounds. The threshold may also depend on \(C_n\).
\end{auxlemma}

\begin{proof}
\proofstep{Point evaluation on a sequence of blocks.}
Use the single compact-finite coefficient representing \(J_\mu\)
in Proposition~\ref{prop:local-operator-normalization}. For any such
coefficient \(F\), denote its transformed squared norm at the point
and its integrated squared norm on the big cell by
\[
 E_F(\gamma)=w(\gamma)\|J_F(\gamma)\|_{\rm HS}^2,\qquad
 {\cal M}(F)=\int_{\text{big cell}} E_F(c)\,dc .
\]
Put \(p=r+1\), \(p_0=2p\), \(N=3p\). In a cut between differing
compact coordinates divide a diagonal word energy by its positive
word weights as in Proposition~\ref{prop:real-node-positivity}.
For the point estimate one needs the three-derivative profile
\[
             {\cal P}(\xi)=\min\{1,(\Lambda/|\xi|)^3\},
             \qquad {\cal P}(0)=1,\qquad \Lambda\le b^{C m}.
\]

A cut of length \(30p^2+6\) admits \(3p(2r+2)=6p^2\) nodes
five units apart and at least two units from the endpoints, since their total required length is \(5(6p^2-1)+4\). In each interpolation take fresh
nodes for the degree (at most \(2r+1\)) of the new letter, using
Proposition~\ref{prop:letter-polynomial-degree}. Use the positive
normalizer \(\prod_i(|Y+\lambda_i|+c_0)\). With \(n_*=30p^2+6\)
and \(A_*=4p+n_*\), one may take
\[
 e=\|\nu\|_1+p(A_*+n_*+N+2),\qquad c_0=2+2e .
\]

Multiplication by the determinant costs \(pA_*\) unit shifts, the
compact contractions cost at most \(pn_*\), and \(N\) radial letters
use at most \(pN\) further shifts. Thus the normalizer includes every
ambient eigenvalue used below. Writing \(h_*\) for the least
unshifted row gap, Lemma~\ref{lem:marked-paired-moment-range} gives
\[
 h_* \ge \min(4\Delta,\eta_{\rm size}/m,b_1)>nm^{10},
 \qquad e\le C_nkm^3+O(m^3),\quad 200 c_0<h_*
\]
on the late family. In particular all shift and degree budgets in
Definition~\ref{hyp:paired-frame-moment-range} allow these endpoints.

For each of the two column polynomials, pair the \(k-1\) interlaced
row roots with the adjacent column roots \emph{away from} the cut.
Each such numerator distance is at most its paired denominator
distance. Omit the two coordinates adjoining the cut if the compact
type has been moved there. Each exceptional numerator is bounded
by \(b^{O(1)}\), and the interior nodes stay a fixed distance from
the column poles and their unit translates.

Padding the row roots by \(c_0\) has a bound even at the actual
endpoint radius \(pq\). Remove a root nearest \(Y\), and bound its
padded factor directly by \(b^{O(1)}\). The \(j\)-th remaining
root on either side has distance at least \(jh_*/2\). Thus the
other ratios have product at most
\[
 \exp\!\left(\frac{4c_0}{h_*}\sum_{j=1}^n\frac1j\right)
        \le (\mathrm{e}n)^{4c_0/h_*}\le\mathrm{e}n .
\]
The other \(m\) ambient factors are each at most \(b^{O(1)}\).
These estimates hold also at the adjacent unit translates of the
nodes. Previous letters change the rational factors there by bounded
ratios for each of the at most \(N\) letters, since distinct nodes
are five units apart. Interpolating a leading coefficient on
integer-spaced nodes has total absolute cost at most \(2^{2r+1}\),
by the product formula for the interpolation weights. This gives
the asserted \(\Lambda\), after incorporating the character and line
norms of Definition~\ref{hyp:paired-frame-moment-range}.

Apply the signed derivative step of
Proposition~\ref{prop:positive-block-fourier} three times, as in
the point recursion following Lemma~\ref{lem:empty-word-exhaustion}. For a size-one block, Fourier inversion of its diagonal
costs at most
\[
              \int_{\C}{\cal P}(\xi)\,d\xi=3\pi\Lambda^2 .
\]
For a larger block choose the coset of the square lattice of side
\(\Lambda^{-1}\) through the prescribed point. Keep that diagonal
in the nonnegative period sum of Lemma~\ref{lem:germ-period}.
The comparison for a block of one smaller size uses the
inverse-coordinate frequencies in Lemma~\ref{lem:gram-limit}.
The \(j\)-th annulus of the dual lattice contains \(O(j+1)\)
frequencies. For a block of size \(s\) this gives
\[
 \sum_{\xi\ {\rm dual}}{\cal P}(\xi)
       \ll 1+\sum_{j\ge1}(j+1)j^{-3}\ll1,\qquad
 E_{h\|(s,t)}\le(C\Lambda^2)^s\max_{|w'|-|w|\le3s}E_{h,w'} .
\]
The power of \(\Lambda^2\) at each unit of size is the reciprocal
covolume of the translation lattice.

Apply this argument first away from \(G=0\) with a smooth cutoff.
At total size \(p\) the cost is \((C\Lambda^2)^p\). At the empty
prefix the limiting normalized energy of every resulting word is
at most \({\cal M}(F)\) when \(F\) is the coefficient multiplied
by \(G^{p_0}\) and the compact contractions below. Indeed the
commuting Casimirs can be moved to this coefficient in the big-cell
pairing by Lemma~\ref{lem:radial-cutoff-transfer}. Tensoring with the
matrix polynomials in \(F\) gives a finite sum of infinitesimal
characters. Different characters are orthogonal under the invariant
pairing, and the positive normalizer bounds each eigenvalue at
every letter. Hence its normalized product has modulus at most one.

The extra matrix polynomial and the cutoff satisfy
\[
 \deg(\text{multiplier})\le 2pp_0+pn_0,\qquad n_0\le30p^2+6,\qquad
                        N=3p\le4p=2p_0 .
\]
The comparison following Lemma~\ref{lem:empty-word-exhaustion}
gives the integrable majorant for both compact types. For a fixed positive
point the cutoff eventually equals one. Removing it gives
\[
                    E_F(\gamma)\le(C\Lambda^2)^p{\cal M}(F).
\]

\proofstep{Obtaining a cut for both compact types.}
For different marked types a differing coordinate gives a cut of
length at least \(30p^2+6\) by
Definition~\ref{hyp:critical-level-weights} and the marked spacings.
Use \(F=G^{p_0}f\) with \(f\) representing \(J_\mu\), and put
\(n_0=0\). When the types agree take \(n_0=30p^2+6\) and choose
a gap in the compact row. The shifts producing the two long cuts are
\[
        u_i-u_{i+1}\ge2n_0+1,\qquad
    u^+=u+n_0\mathbf1_{\{i+1\}},\quad u^-=u-n_0\mathbf1_{\{i\}} .
\]

For clarity, let \(L_1>\dots>L_k\) be the strictly decreasing integer
coordinates of that compact row after subtracting its common scalar.
On restriction to the stabilizer of a unit vector the type labels
\(M\) satisfy \(L_{j+1}<M_j\le L_j\). The squared projections when
one upper coordinate \(L_a\) is raised or lowered are respectively
\[
 \frac{\prod_{j<k}|M_j-1-L_a|}{\prod_{j\ne a}|L_j-L_a|},
 \qquad
 \frac{\prod_{j<k}|M_j-L_a|}{\prod_{j\ne a}|L_j-L_a|}.
\]
These are the residues of the compact branching polynomials in
Lemma~\ref{lem:residual-single-cofactors}: taking the \(a\)-th
spectral projection divides the characteristic polynomial of the
compressed operator by the derivative of the full polynomial at
\(L_a\). The raising polynomial has roots \(M_j-1\) and the lowering
polynomial roots \(M_j\), from the unit shift in the dual action.

For \(M_i\ge L_{i+1}+n_0+1\) the entire raising chain is allowed.
Otherwise \(M_i\le L_i-n_0\) and the lowering chain is allowed.
In any allowed step write \(b_j=\max(1,|L_j-L_a|)\).
The \(j\)-th numerator factor is at least \(b_{j+1}\) for \(j<a\)
and \(b_j\) for \(j\ge a\), by interlacing and integrality.
Multiplying the lower bounds gives the cancellation
\[
 \frac{\prod_{j<a}b_{j+1}\prod_{a\le j<k}b_j}
      {\prod_{j\ne a}b_j}
       =\frac{1}{b_1b_k},\qquad b_a=1 .
\]
If \(B_1=b^{C'}\) bounds the compact entries
and the shifted width is at most \(3B_1\), each squared projection
on the allowed chain is at least \((3B_1)^{-2}\).
Both chains have operator norm at most one at a unit vector.

To obtain the contraction vector at \(\gamma\), take the last row of the
adjugate of the cross-pairing matrix. Combine the left frame with
its conjugate coefficients, transport by \(\gamma^{-1}\), then
project onto the right positive middle space perpendicular to \(z'\).
The adjugate identity annihilates its pairings with the first \(r\)
right columns. Its pairing with \(z'\) has modulus \(|\Delta_p|\).
Before subtracting its \(z'\)-component only the left distinguished
line contributes to the norm. Its squared positive length is thus
\[
               |\Delta_p|^2/T'-T|\Delta_r|^2=G/T' .
\]

Let \(A_\pm\) be the two contraction chains evaluated at this vector,
oriented from the shifted compact type to the original one. Sum the
squared lower bound over the orthogonal \(M\)-components and use
homogeneity of degree \(n_0\). For any \(y\) in the original type,
\[
 \|A_\pm\|^2\le(G/T')^{n_0},\qquad
 \sum_\pm\|A_\pm^*y\|^2\ge (3B_1)^{-2n_0}(G/T')^{n_0}\|y\|^2 .
\]
The contraction is invariant under both unipotents and the left
stabilizer. Right compact translation changes \(A_\pm\) by its
two compact actions, which cancel those in the orbital integral.
Thus \(J_{F_\pm}=G^{p_0}J_f A_\pm\) for \(F_\pm=G^{p_0}f A_\pm\).
The same point, stabilizer Haar and quotient measure occur in all
three integrals. Taking adjoints and summing over a target
orthonormal basis gives
\[
 E_f(\gamma)\le(3B_1)^{2n_0}(T')^{n_0}G^{-A}
                 \sum_\pm E_{F_\pm}(\gamma),\qquad A=2p_0+n_0.
\]
For different types the same estimate has just one term and \(A=2p_0\).

\proofstep{The determinant moment in the integrated norm.}
On the big cell \(0<G\le|\Delta_p|^2\). Hence
\[
 {\cal M}(F_\pm)\le (T')^{-n_0}
       \int_{\text{big cell}}|\Delta_p|^{2A}E_f(c)\,dc .
\]
Use \eqref{eq:minor-shift-unit} with \(l=p\), now \(A\) times,
in the Hilbert space of square-integrable operator-valued sections.
Each nonzero endpoint has norm one. Minkowski's inequality gives
\[
 \|\Delta_p^A U_\mu\|_{\rm cell}\le B^A,\qquad
 \int_{\text{big cell}}|\Delta_p|^{2A}E_f(c)\,dc\le B^{2A}M_\mu .
\]
Here \(A\le40p^2\), so the extra radius \(Ap\le40p^3\) fits
the window of Lemma~\ref{lem:minor-orbital-shift} after increasing
\(C_n\) by one. This use of its finite expansion precedes any
bound on a point value. The factors \((T')^{n_0}\) cancel.
The two remaining factors satisfy
\[
    2(C\Lambda^2)^p(3B_1)^{2n_0}B^{2A}\le b^{O(m^4)},\qquad
                   G^{-A}\le\max(1,G^{-1})^{10m^2}.
\]

Finally, transporting the character to unit size by the flag torus
scales its \(j\)-th column by \(\alpha^j\) for \(0\le j\le r\).
The line and Gram scalings are
\[
 (T,T')\longmapsto\alpha^{2r}(T,T'),\qquad
    G=|\Delta_p|^2-TT'|\Delta_r|^2\longmapsto\alpha^{2rp}G .
\]

For a prefix of length \(a\) the minor scales by
\(\alpha^{a(a-1)}\). In each ordinary block the \(2k\) part of the
density exponent multiplies the ratio of the two block-end minors.
Across the blocks these terms telescope. If the last block is a
singleton the identity \(G=|\Delta_r|^2(|t_0|^2-TT')\) gives the
same total exponent \(2krp\) as on the big cell. The common
\(k\)-dependent dilation cancels in the ratio of the two energies.

In big-cell pivot coordinates the \(j\)-th complex coordinate
scales by \(\alpha^{2j}\), so its section Jacobian is
\[
             J_{\rm sec}=\prod_{j=0}^r\alpha^{4j}=\alpha^{2rp}.
\]

In a block of size \(s\) the within-block root coordinates are
indexed by pairs \(a<b\). Write \(D_{\rm rt}=\sum(b-a)\) over these pairs and
\(c=\sum_{\rm blocks}\binom{s}{2}\). On a block starting after
\(j\) columns the remaining density exponent, relative to splitting
the block into singletons, is
\[
 4\sum_{i=0}^{s-1}(j+i)(2i+1-s)
       =4\sum_{0\le a<b<s}(b-a).
\]

The root \(ab\) has complex-coordinate Jacobian
\((\alpha^b/\alpha^a)^2\) and contributes \(\alpha^{-1}\)
in the unit-root normalization. Thus the square of their
combined orbital change has exponent \(4D_{\rm rt}-2c\).
Canceling against the density and section changes leaves the
factor \(J_{\rm sec}\alpha^{-2c}\), or its reciprocal on reversal. Both therefore satisfy
\[
 \max\{J_{\rm sec}\alpha^{-2c},(J_{\rm sec}\alpha^{-2c})^{-1}\}
   \le\max(\alpha,\alpha^{-1})^{2rp+2c},\qquad 2c\le rp .
\]

Multiplication by \(G^{-A}\) costs at most a further \(2rpA\)
powers of this scale. The independent root Haar factor \(\theta'\)
in \eqref{eq:absolute-root-scales} satisfies
\(b^{-|C_1|m^2}\le\theta'\le b^{|C_1|m^2}\), so its inverse square
costs at most \(b^{2|C_1|m^2}\). Since
\(\max(\alpha,\alpha^{-1})\le b^{O(1)}\), all these powers are
absorbed in \(b^{C_bm^4}\).
\end{proof}


The bounds on products of height losses can now be inserted into the local norm estimates. The resulting inequalities concern a single normalized orbital term and are ready to be multiplied for a flag pair.

\begin{proposition}[Two bounds for a single amplitude]\label{prop:two-row-amplitude-bounds}
Use the normalized local Bessel bounds of
Sections~\ref{sec:local-amplitudes} and~\ref{sec:absolute},
the height hypotheses of Lemma~\ref{lem:row-height-products}, and
the minor iteration of Lemma~\ref{lem:minor-orbital-iterate}
with \(pq\le C_{\rm pr}km^3\).
On the selected range where
\(T'/T\le A\), \(T_{{\rm crit},v}g^{-m}\ge k^{3r}\),
\(r+2\le\delta_m k\), \(C_{\rm pt}\le C_a\le k\),
\(Cm\ge2r+4\), \(e\pi\le k\), and \(C_a\ge6e\pi A\), one has
\begin{equation}\label{eq:row-l2}
 n_Z A_t^2\le L_Z,\qquad
 L_Z=H_Z^{2h}\Gamma_S\tau\,\sigma_Z^{C_2m^2-2a}.
\end{equation}
The additional estimate needed when a pair has a degenerate determinant is
\begin{equation}\label{eq:row-abs}
 n_Z A_t^2 K_Z^{2k}T^{-kf}\le
 H_Z^{2h}\Gamma_S E^2(C_a/k)^{kf}
                 \sigma_Z^{C_2m^2+2C_2m^4-2a}.
\end{equation}
\end{proposition}

For the two marked words Proposition~\ref{prop:assigned-prime} gives \(A=2\) in either orientation. Take \(C_a=\max(1,C_{\rm pt},12e\pi)\) and \(k\ge C_a^3\) at the common threshold. 
We prove both bounds in the next subsection. The second estimate
will cancel the factor \(T^{-kf}K_X^kK_Y^k\) in a pair count.

\subsection{Local conversion from the operator norm}

The proposition concerns normalized row terms, whereas the local estimates concern orbital operators. The conversion uses the cell density at every block end as well as the dimension of the definite type.

\begin{proof}[Proof of Proposition~\ref{prop:two-row-amplitude-bounds}]
\proofstep{Gram factor and cell density.}
The cell density \(w\) is related to the Gram factor by the
block determinant. At one place put \(d_a=|\Delta_a|\) at
each block end and \(E(a)=4(p-a)+2k\). An ordinary block of
size \(s\) from \(a\) to \(a+s\) contributes
\((d_{a+s}/d_a)^F\), where
\(F=2(2(p-a-s)+k+s-2)\).
The block contributes the ratio
\[
 \frac{d_{a+s}^{E(a+s)}}{(d_{a+s}/d_a)^F d_a^{E(a)}}=
 d_{a+s}^{E(a+s)-F}d_a^{F-E(a)}
     =d_{a+s}^{4-2s}d_a^{-(2s+4)} .
\]
Both exponents have absolute value at most \(m+4\). Start with
\(G_{Z,v}^k\le d_p^{2k}\) and insert the blocks backwards.

For a terminal singleton with pivot \(t_0\) the factor is
\((|t_0|^2-TT')^{k-1}\).
In that case \(G_{Z,v}=d_r^2(|t_0|^2-TT')\), so the
initial ratio is \(G_{Z,v}d_r^{-6}\).
Since \(d_0=1\), the resulting comparison in both cases is
\[
 G_{Z,v}^k\le w\,\max(G_{Z,v},G_{Z,v}^{-1})
       \prod_a\max(d_a,d_a^{-1})^{2(m+4)} .
\]

\proofstep{Normalization of the local squared norm.}
Let \(M_\lambda\) be the squared norm of the local amplitude
in the common Bessel measure, and set \(t=T_{\rm crit}/T'\).
For the squared-norm estimate one first bounds the point value
of the unit-cell-norm amplitude \(U_\mu=J_\mu/\sqrt{M_\mu}\).
Write \(w\) for the cell density computed above.
Lemma~\ref{lem:unit-cell-evaluation} bounds
\(w\|U_\mu(\gamma)\|_{\rm HS}^2\) on the full endpoint
window \(\|\mu-\lambda\|_1\le pq\).
Proposition~\ref{prop:local-operator-normalization} fixes \(M_\mu\).
The formal-degree
and compact Weyl normalizations give
\[
 (e\pi g^{-m}/(k\upsilon T))^k M_\lambda
       \le d_v\,t^{-k}b^{C_c m^2}
\]
at the base parameter.

To compare the constants, write
\(\mathcal F(j)=\prod_{i=1}^{j-1}i!\) and use
\(\operatorname{vol}U(j)=(2\pi)^{j(j+1)/2}/\mathcal F(j)\).
Expansion of the volume quotient cancels the compact
dimension denominator \(\mathcal F(k)\) and leaves the
factor \((e/k)^k(k-1)!\le e\). The remaining level and volume
factor is bounded by
\[
 (2\upsilon g)^{(r+1)(2r+1)}
       \le b^{(1+2C_g)m^2/2}
\]
using \(2\upsilon\le b\), \(g\le b^{2C_g}\) and
\((r+1)(2r+1)\le m^2/2\). The Schur comparison costs at most
\(b^{C_Sm^2}\). This gives the stated
bound with \(C_c\ge C_S+1+(1+2C_g)/2\).
The cell-height factor and the shifted pointwise loss
combine with \(b^{C_c m^2}\) in \(L_v\).

At one place let \(D=|\Delta_r|\),
\(n=(e\pi g^{-m}G_{Z,v}/(k\upsilon T))^k\), and
\(s_0=s_v^{-q/p}\). The \(L^2\) bound has the form
\(n\|J\|_{\rm HS}^2/d_v\le L_v\,t^{-k}s_0^2\).
For \(D>0\) the pointwise bound is
\[
 \|J\|_{\rm op}\le
   P_v\,(C_{\rm pt}/(k\upsilon))^{r(k-r-2)}D^{-(k-r-2)}s_0.
\]
For the Hilbert--Schmidt inequality, apply
Lemma~\ref{lem:minor-orbital-iterate} to \(U\) with the square
root of \eqref{eq:unweighted-unit-cell-evaluation} as endpoint
bound. Multiply back by \(\sqrt{M_\lambda}\) only at the base.
For the operator inequality take \(V=J\) and apply
Proposition~\ref{prop:absolute-cell-operator} at every nonzero
endpoint. Its \(D,D_e\), measures and constants are the same
there. Thus the two local multipliers are respectively
\[
   s_v^{-2q/p}\quad\hbox{in the squared estimate},\qquad
   s_v^{-q/p}\quad\hbox{in the operator estimate}.
\]
The second application requires \(D>0\); the first uses
\(G_{Z,v}>0\) and permits \(D=0\).

For the first global inequality multiply these local squared-norm
bounds. Tensor factorization supplies \(n_Z A_t^2\) with a factor
of the term volume squared. On the right one obtains
\[
 H_Z^{2h}\left(\prod_v L_v\right)
          \left(\prod_v t_v^{-k}\right)\sigma_Z^{-2a}.
\]
Lemma~\ref{lem:row-height-products} and the word-weight estimate give \eqref{eq:row-l2}.

\proofstep{A local gain from the preceding minor.}

The squared norm applies on the whole allowed cell. An additional gain becomes available when the preceding minor is nonzero, the hypothesis needed for the absolute orbital estimate.

The rank estimate \(\|J\|_{\rm HS}^2\le d_v\|J\|_{\rm op}^2\)
turns the pointwise bound into a bound for the weighted square.
For \(D>0\) the preceding minor is a block end, and
\(D^{2r+4}\le L_v\) since \(Cm\ge2r+4\).
After using this inequality the squared pointwise bound
with weight \(n\) is at most
\[
 L_v\,P_v^2s_0^2
 \left(\nu g^{-m}(C_{\rm pt}/k)^{2r(1-(r+2)/k)}\right)^k,
 \qquad \nu=G_{Z,v}/(TD^2).
\]
Here \(\upsilon\ge1\) and \(e\pi\le k\).
Take the minimum with the \(L^2\) estimate. For \(C_{\rm pt}\le C_a\le k\)
and \(r+2\le\delta_m k\), the resulting multiplier after
factoring out \(t^{-k}\) is bounded by
\[
 \min\bigl(1,\nu t g^{-m}(C_a/k)^{2r(1-\delta_m)}\bigr)^k.
\]
\proofstep{Combining the two estimates at each place.}
The case \(D=0\) uses the \(L^2\) bound with cap equal to one.

For \(\mathrm{sh}\ge1\) the local Cayley factor satisfies
\(K_v^2\le(1+\sqrt2)^2<6\). The critical-level estimate
\(T'tg^{-m}\ge k^{3r}\) then gives
\[
 \frac{n\|J\|_{\rm HS}^2K_v^{2k}}{d_v T^k}
 \le L_v\,s_0^2(6A g^{-m}k^{-3r})^k,
 \qquad T'/T\le A .
\]
The case \(D=0\) uses the same bound with \(K_v=1\).
For \(0<\mathrm{sh}<1\) insert \(G_{Z,v}=TT'D^2\mathrm{sh}^2\)
in the squared operator-norm bound. The factors
\(\mathrm{sh}^{2k}\) cancel, and
\((1+\sqrt{1+\mathrm{sh}^2})^2<6\).
The pivot loss absorbs \(D^{2r+4}\) as before, leaving
\(L_v\,P_v^2s_0^2(6e\pi A/k)^k\).
Taking the absolute bound constant at least \(6e\pi A\)
gives \eqref{eq:row-abs} upon taking products, using the
place-product bounds for the losses.
\end{proof}

\subsection{The sums at corank at least two}

At higher corank two independent sums remain, over the Gram labels and over the term pairs with given labels. Decay in the label height will make the first sum finite after the fibre count has been inserted.

\begin{lemma}[Summation over reduced Gram labels]\label{lem:gram-label-sum}
For \(\phi(Z)=\sigma_Z^{-12p^2}\) the sum over the reduced labels in
the row satisfies 
\[
 \sum_Z\phi(Z)\le\Phi=(C_*^fP^3)^{p^2},
\]

with the absolute constant \(C_*\) of the lattice packing bound.
\end{lemma}

\begin{proof}
Each of the \(p^2\) scaled entries lies in a fixed lattice
of multiplicative minimum one. The sum of its weight
\(\prod_v\max(1,|z_v|)^{-3}\) is at most \(C_*^f\) by the
packing estimate, with \(C_*\) absolute. The entry height
is at most \(P\sigma_Z^4\). The product of the entry weights
therefore dominates \(\phi(Z)/P^{3p^2}\).
The entries determine \(Z\), yielding
\(\sum_Z\phi(Z)\le\Phi=(C_*^fP^3)^{p^2}\).
\end{proof}

\begin{lemma}[Summing estimates on a label fibre]\label{lem:row-fibre-weight}
Let \(M\) be the row mass and \(N(X,Y)\) a bound for the
weighted size of the fibre divided by \(M\). Suppose its two
amplitudes divided by \(\sqrt d\) have bounds \(\alpha(X)\)
and \(\beta(Y)\). If
\(N(X,Y)\alpha(X)\beta(Y)\le R\phi(X)\phi(Y)\),
its contribution \(C(X,Y)\) satisfies 
\[
 |C(X,Y)|\le MdR\phi(X)\phi(Y).
\]

\end{lemma}
\begin{proof}
Write \(w_{t,t'}\ge0\) for the counting weight of an ordered
pair in the fibre. Then
\[
 \sum_{(t,t')\mapsto(X,Y)} w_{t,t'}|a_t a_{t'}|
 \le d\alpha(X)\beta(Y)\sum_{(t,t')\mapsto(X,Y)}w_{t,t'}
 \le MdN(X,Y)\alpha(X)\beta(Y).
\]
The same estimate holds for any finite partial sum of the fibre,
so the weighted count also bounds the full absolute sum.
\end{proof}

Summability in each label leaves the number of pairs in a fixed fibre. At corank at least two this count depends on whether the joint orthogonal determinant degenerates.

\begin{proposition}[Pairs of corank at least two]\label{prop:row-higher-corank}
For \(s\ge2\) in the pair count set
\(E_s=c((s+m)m+2+\log_2 f)\). In the nondegenerate case
\[
 a\ge4ck+4E_k+4h+|C_2|m^2/2+12p^2
\]
suffices to leave the summable weight \(\phi(X)\phi(Y)\) after
applying \eqref{eq:row-l2}. In the degenerate case adding
\(|C_2|m^4\) to this lower bound suffices. The two amplitudes
together gain \((C_a/k)^{kf}\) from \eqref{eq:row-abs}.
More precisely, let \(S_s^{\rm nd}\) and \(S_s^{\rm deg}\) be the
absolute pair sums in these two cases, in units of the total row
mass times \(d=\dim\pi_u\), before the finite frame factor is inserted.
For \(2\le s\le k\) put
\[
 W_s=\Phi^2\Gamma_S (csf)^{6sf}P^{2(E_s+h)}\Xi_s(P^2),
\]
where \(\Xi_s\) is the factor in Proposition~\ref{prop:flag-pair-counts}.
Then
\begin{equation}\label{eq:high-corank-summed-prices}
 \begin{aligned}
 S_s^{\rm nd}
 &\le W_s\,\tau\,(2/g)^{(s-1)kf}(k\upsilon/(e\pi))^{kf},\\
 S_s^{\rm deg}
 &\le W_s\,E^2(2/g)^{(s-2)kf}(C_a\upsilon/(e\pi))^{kf}.
 \end{aligned}
\end{equation}
\end{proposition}
\begin{proof}
For corank \(s\ge2\) the count has factors
\[
 \frac{D_{\rm pair}^k}{N_T^k}\,
 \Xi_s(H_XH_Y)(H_XH_Y)^{E_s}(csf)^{6sf},
 \qquad E_s=c((s+m)m+2+\log_2 f).
\]
By \eqref{eq:flag-pair-mass-factor}, \(\Xi_s(H_XH_Y)\) is constant in the
heights for \(s\le k-2\). At \(s=k-1,k\) it contributes the
additional factor \((H_XH_Y/P^2)^{ck}\) relative to its value at \(P^2\).

In the nondegenerate case use the bound
\(D_{\rm pair}\le(2/g)^{(s-1)f}\sqrt{G_XG_Y}\).
The product of the two normalization factors is
\[
 (n_Xn_Y)^{-1/2}
    =\left(\frac{k\upsilon}{e\pi}\right)^{kf}
        \frac{N_T^k}{(G_X G_Y)^{k/2}}.
\]
It cancels the \(N_T^{-k}\) and the Gram powers in the
determinant estimate, leaving the factor
\[
 (2/g)^{(s-1)kf}\left(\frac{k\upsilon}{e\pi}\right)^{kf}.
\]

Set \(j_s=0\) for \(2\le s\le k-2\) and \(j_s=ck\) otherwise.
For \(Z=X,Y\), the variable part of the height bound is
\[
 (H_Z/P)^{j_s+E_s+h}\le \sigma_Z^{4j_s+4E_s+4h}.
\]
The constant part from both labels is exactly bounded by
\(\Xi_s(P^2)P^{2(E_s+h)}\). After the nondegenerate amplitude loss the net exponent is bounded by
\[
       4ck+4E_s+4h+C_2m^2/2-a\le-12p^2 .
\]
The last inequality uses the assumed lower bound on \(a\),
so Lemma~\ref{lem:row-fibre-weight} applies with the indicated
summable weight.

For the degenerate case the determinant bound uses
\((2/g)^{(s-2)f}\sqrt{G_XG_Y}T^{-f}K_XK_Y\).
Indeed \eqref{eq:row-abs} bounds
\(\sqrt{n_Z}|A_t|K_Z^k T^{-kf/2}\) by the square root
of its right-hand side. Multiply the two square-root bounds. Their extra factors become
\[
      (K_X^kT^{-kf/2})(K_Y^kT^{-kf/2})
                 =T^{-kf}K_X^kK_Y^k,
\]
exactly the additional factor in \(D_{\rm pair}^k\). The two determinant
bounds lie below \((H_XH_Y)^3\) by the bound on \(K_Z\)
and the height allowances \(G_Z\le H_Z\) and \(3^f\le H_Z\).

In a nonempty fibre the integral determinant is at least one.
The degenerate row gains \((C_a/k)^{kf}\) and costs an
additional \(|C_2|m^4\) in the required lower bound on \(a\).
Thus the indicated choices of \(a\) leave the same exponent
\(-12p^2\) on each label. Summing the two weights costs one
\(\Phi^2\) by Lemma~\ref{lem:gram-label-sum}. The two square
roots of the pointwise bounds contribute \(\Gamma_S E^2\);
those of the squared-norm bounds contribute \(\Gamma_S\tau\).
This gives both rows of \eqref{eq:high-corank-summed-prices}.
\end{proof}

\begin{auxlemma}[The total error from pairs of corank at least two]\label{lem:higher-corank-total-error}
For the marked words \(u,u'\) of Proposition~\ref{prop:assigned-prime},
let \(Q=kf\log k\), \(d=\dim\pi_u\), and
\(\theta_{uu'}=d_ud_{u'}/D_*\). Take the moment order of
Lemma~\ref{lem:row-moment-budget} in \eqref{eq:high-corank-summed-prices}.
Let \(M_u\) be the total row mass for \(T=T_u\), and let
\(\Gamma_1\) be the finite diagonal squared norm of a coarse frame at the
two split type primes. For sufficiently large \(x\), uniformly in
\(u,u'\),
\begin{equation}\label{eq:high-corank-in-gram-units}
 \Gamma_1 M_u d\sum_{s=2}^k(S_s^{\rm nd}+S_s^{\rm deg})
       \le e^{-Q/2}\theta_{uu'} .
\end{equation}
Both terms in every pair here have positive single-term Gram invariant.
\end{auxlemma}
\begin{proof}
\proofstep{Decay of the middle-corank mass.}
Write \(L=\log b\) and \(\upsilon=\operatorname{rd}(E)/2\). In
\(2\le s\le k-2\) put \(a_s=\min(s,k-s)\), \(t=k-a_s\) and
\(n=a_st=s(k-s)\). With \(R_i=m_i/m_{i-1}\), cancellation gives
\[
 \Xi_s=\frac{m_1}{R_{t+1}}\prod_{i=2}^{a_s}\frac{R_i}{R_{t+i}} .
\]

In \eqref{eq:adjacent-unimodular-genus-masses}, each Euler product
at \(i\ge2\) lies between \(\zeta(i)^{-f}\) and \(\zeta(i)^f\),
hence between \(2^{-f}\) and \(2^f\).
The factorial denominator from \(R_{t+1}\) is \(t!\). In each
remaining quotient it is \((t+i-1)!/(i-1)!\ge t!\).
The power of \(\upsilon/\pi\) in the denominator is
\(t+1+(a_s-1)t=n+1\) at each real place. Therefore
\begin{equation}\label{eq:middle-corank-factorial}
 \Xi_s\le m_1(2\upsilon)^{f/2}4^{a_sf}
               (\pi/\upsilon)^{(n+1)f}(t!)^{-a_sf}.
\end{equation}

In particular \(\Xi_2=m_2/(R_{k-1}R_k)\). The degenerate level
factor at \(s=2\) is exactly one. Multiplication by
\((C_a\upsilon/(e\pi))^{kf}\) still leaves the power
\(\upsilon^{-(k-3)f}\) in \eqref{eq:middle-corank-factorial},
apart from the rank-one factor \(m_1(2\upsilon)^{f/2}\).

For both rows of \eqref{eq:high-corank-summed-prices}, bound the
archimedean factor by \((Ck\upsilon)^{kf}\).
Since \(n\ge2(k-2)\), \(t\ge k/2\) and \(\upsilon\ge1\), the volume
power \(\upsilon^{(k-n-1)f}\) is at most one for \(k\ge4\).
By the rank-one and class-number estimates of
Section~\ref{sec:flag-counts},
\[
 m_1\le2^f(2\upsilon)^{f(1+\log_2(2f))},\qquad
 \log\!\big(m_1(2\upsilon)^{f/2}\big)\ll f(1+\log f)L .
\]
Now \(t!\ge(t/e)^t\). As \(f\le k\), \(2\upsilon\le b\) and
\(L=o(k)\) on the chosen row, \eqref{eq:middle-corank-factorial} gives
\[
 \log\!\big(\Xi_s(Ck\upsilon)^{kf}\big)
 \le -(n-k)f\log k+O(nf)+O(f(1+\log f)L)
 \le -(1-o(1))Q .
\]
For large \(k\) the first two terms together decrease
with \(n\), so inserting \(n=2(k-2)\) gives a uniform estimate.

\proofstep{Summation of the heights and the extreme coranks.}
From the definitions of \(P,\Phi,\Gamma_S,E,E_s,h\),
the logarithm of the height, packing and label factors in
\(W_s E^2/\Xi_s(P^2)\) is at most
\[
 C\{(s+m)m^4 fL+m^6fL+sf\log(Csf)\}.
\]
Here \(2+\log_2 f\le m\). The selected size bounds
\(m\le x,\ k\ge m^5Lx^2\) imply \(m^6L=o(k)\).
For \(s=2\) the displayed cost is \(o(Q)\).
For \(s\ge3\), its \(s\)-dependent terms are
\(o((s-2)kf\log(g/2))\), since \(s/(s-2)\le3\),
\(sf\le k^2\), and \(g_i\ge(kb)^8\). The remaining terms are \(o(Q)\).
Thus even the smaller level exponent \(s-2\) pays the
\(s\)-dependent loss in every middle corank \(s\ge3\).

For \(s=k-1,k\) the rank-zero or rank-one packing estimate uses
\(\Xi_s(P^2)=(ck)^{kf}(P^2 N_T)^{ck}\).
By Proposition~\ref{prop:assigned-prime}, \(\log T\ll mL\), so
\[
 \log\Xi_s(P^2)\ll kf\log k+k m^3 fL .
\]
The archimedean logarithm is \(O(kf(\log k+L))\).
Both are \(o((k-3)kf\log(g/2))\), using \(m^3=o(k)\) and
\(\log(g/2)\gg L+\log k\). Together with the preceding height bound
this leaves an exponent \(\le-2Q\) in either extreme corank,
before restoring the word dimensions.

\proofstep{Conversion to the Gram normalization.}
Proposition~\ref{prop:flag-row-mass} gives \(M_u=m_k(T_u^k-1)^f\).
At a split type prime of residue cardinality \(q\) the diagonal
finite norm is \(|\mathrm{GL}_k(\mathbb F_q)|/|U_k(\mathbb F_q)|\)
by Lemma~\ref{lem:coarse-split-gram}, where \(U_k(\mathbb F_q)\)
is the upper unitriangular subgroup. Its product over both primes
and the \(f\) places is \(I_k\) of Proposition~\ref{prop:marked-column-dimension}.
The definite dimension there is precisely \(d\), hence
\[
 \Gamma_1 M_u d=I_k\,m_k(T_u^k-1)^f d=d_u .
\]

In Gram units the nondegenerate row has
\(d_u\tau/\theta_{uu'}=1+\delta\). For the degenerate row, use
the lower endpoint \(T_{u',*}\) of the prime interval in Proposition~\ref{prop:assigned-prime}.
The defining equation for the target, \(T_{u'}\ge T_{u',*}\)
and Bernoulli's inequality give the lower end of \eqref{eq:linked-rank-window},
\[
 \frac{d_{u'}}{D_*}
 =\Gamma_S^{-2}(T_{u'}/T_{u',*})^{kf}(1-T_{u'}^{-k})^f
 \ge \Gamma_S^{-2}(1-f2^{-k})\ge(2\Gamma_S^2)^{-1}.
\]
Thus the logarithmic cost is at most \(\log(2\Gamma_S^2)=o(Q)\).
After these conversions each of the two bounds at each corank
is at most \(e^{-3Q/4}\theta_{uu'}\). Finally \(2k\le e^{Q/4}\)
for large \(k\), giving \eqref{eq:high-corank-in-gram-units}.
\end{proof}

\begin{construction}
\begin{tabular}{@{}lll@{}}\toprule
Pair determinant & Single-term estimate & Level factor in the pair sum\\\midrule
Nondegenerate & \eqref{eq:row-l2} & \((2/g)^{(s-1)kf}\)\\
Degenerate & \eqref{eq:row-abs} & \((2/g)^{(s-2)kf}\)\\\midrule
Corank & \multicolumn{2}{l}{Mass or packing factor for both determinant cases}\\\midrule
\(s=2\) & \multicolumn{2}{l}{\(\Xi_2=m_2/(R_{k-1}R_k)\)}\\
\(3\le s\le k-2\) & \multicolumn{2}{l}{\(\Xi_s=m_s m_{k-s}/m_k\)}\\
\(s=k-1,k\) & \multicolumn{2}{l}{\(\Xi_s(P^2)=(ck)^{kf}(P^2N_T)^{ck}\)}\\\bottomrule
\end{tabular}
\caption{Factors in the high-corank counts. For a degenerate pair at \(s=2\) the level factor equals one. The amplitude gain \((C_a/k)^{kf}\) leaves \((C_a\upsilon/(e\pi))^{kf}\) after normalization. The two adjacent-rank mass denominators supply the factorial decay in \eqref{eq:middle-corank-factorial}.}\label{fig:high-corank-amplitude-choice}
\end{construction}

\subsection{Corank one}

For the sharp estimate between different words, a nonscalar pair
of corank one forces a gain on at least one side. The distinction
records what happens to the distinguished negative vectors when the
two terms are compared.

\begin{definition}[The scalar relation of a pair]\label{def:corank-one-scalar}
Write \(F,F_j\) for the two flags of term \(j=1,2\), so that \(F\)
is common. Write \(X,X_j\) for their isotropic spans and \(z,z_j\)
for their distinguished vectors. Put
\[
 V_j=X+Ez+X_j+Ez_j,\qquad
 Y_j=V_j\cap(X+Ez)^\perp,\quad Y'_j=V_j\cap(X_j+Ez_j)^\perp .
\]
For \(y\in Y_j,\ y'\in Y'_j\) record the relation
\[
 {\cal R}_j(y,y')=
 \{(\alpha,\alpha')\in E^2:
 y-y'-\alpha z-\alpha'z_j\in X+X_j\}.
\]
Let \(\beta\) be the comparison isometry taking \(F_1\) to \(F_2\).
A pair of common rank \(k-1\) is \emph{scalar} if for all
\(y_j\in Y_j,\ y'_j\in Y'_j\) one has
\[
 y_1-y_2\in X,\quad y'_2-\beta y'_1\in X_2
       \ \Longrightarrow\
 {\cal R}_1(y_1,y'_1)={\cal R}_2(y_2,y'_2).
\]
\end{definition}

\begin{lemma}[A small term in a corank-one pair]\label{lem:row-corank-one-gain}
For a nonscalar pair of generic terms of corank one assume
\(0\le\delta_m\le1/4\), \(\max(1,C_a)^3\le k\), and
\(\prod_v t_v^k\le2C_0\Gamma_S^2\). At least one of the two terms
satisfies
\[
 n_Z A_t^2\le L_Z E^2\sigma_Z^{2C_2m^4}Q,\qquad
 Q=2C_0\Gamma_S^2
       (4\sqrt{T'/T}\,g^{-(2r+1)}k^{-r})^{kf}.
\]
\end{lemma}

\begin{proof}
For the nonscalar pair of two generic terms, the corank-one
gluing alternative gives a product lower bound
\((\sqrt{TT'}/g)^f\) for the mixed factors. When both preceding
minors are nonzero, at each place the mixed factor is
\[
 \frac{\sqrt{1+x^2}+\sqrt{1+y^2}}{xy}
 \ \le\ \frac4{\min(1,x)\min(1,y)},\qquad
 x=\mathrm{sh}_v(X),\quad y=\mathrm{sh}_v(Y).
\]
Write \(P_X=\prod_v\min(1,\mathrm{sh}_v(X))\) and use \(P_Y\)
for the other term. Multiplying the mixed-factor estimate yields
\[
      \min(P_X^2,P_Y^2)\le P_XP_Y
                          \le (4g/\sqrt{TT'})^f .
\] When one preceding minor vanishes,
the mixed factor is \(1/\mathrm{sh}\) for the other term.
The same bound then follows for the other term since its
truncated product is at most one. Two zero preceding minors
would contradict the positive mixed product. It follows that
one term, with \(\nu_v=T'\mathrm{sh}_v^2\), satisfies
\(\prod_v\min(\nu_v,T')\le(4g\sqrt{T'/T})^f\).

On that term put \(t_v=T_{\rm crit,v}/T'\),
\(C'=\max(1,C_a)\) and \(u=(C'/k)^{2r(1-\delta_m)}\).
With \(0\le\delta_m\le1/4\) and \((C')^3\le k\) we have
\(k^{-2r}\le u\le k^{-r}\). The critical-level estimate
gives \(T't_v g^{-m}\ge k^{3r}\). Consequently
\[
 \min(1,\nu_vt_v g^{-m}u)
       \le \min(\nu_v,T')t_v g^{-m}u .
\]
The inequality follows directly for \(\nu_v\le T'\), and the
right side is at least one for \(\nu_v>T'\).
Multiplying the surviving local factors gives the stated \(Q\)
after taking the \(k\)-th power and using the bound on \(\prod_v t_v^k\).
\end{proof}

The improvement just proved can occur on either side of a pair. To sum it over a fibre, the count must work in both orientations.

\begin{proposition}[Summed contribution at corank one]\label{prop:summed-row-corank-one}
For the nonscalar generic pairs of Lemma~\ref{lem:row-corank-one-gain},
if \(a\ge12h+|C_2|m^2/2+|C_2|m^4+12p^2\), let \(C_1(X,Y)\) denote their contribution in the fibre of Lemma~\ref{lem:row-fibre-weight}. In units of the row mass \(M\) and dimension \(d\),
\[
       \frac{1}{Md}\sum_{X,Y}|C_1(X,Y)|\le 2c^fP^{6h}\Gamma_S E\tau\sqrt Q\,\Phi^2,
\]
where \(\Phi\) is the bound of Lemma~\ref{lem:gram-label-sum}.
\end{proposition}
\begin{proof}
In the fibre count, the switch and partner estimates in the
two orientations give
\[
 N(X,Y)\le c^f(H_XH_Y)^h\min(H_X^h n_X,H_Y^h n_Y).
\]
The two amplitudes contribute the denominator \(\sqrt{n_Xn_Y}\).
Writing the counting estimate in those units gives
\[
 \frac{\min(H_X^h n_X,H_Y^h n_Y)}{\sqrt{n_Xn_Y}}
              \le (H_XH_Y)^{h/2}.
\]

Use Lemma~\ref{lem:row-corank-one-gain} on one side and
\eqref{eq:row-l2} on the other. For
\(a\ge12h+|C_2|m^2/2+|C_2|m^4+12p^2\) the remaining factors
are bounded by \(c^fP^{6h}\Gamma_S E\tau\sqrt Q\,\phi(X)\phi(Y)\).
Indeed the three height powers on each label are \(h\) from the
count, \(h/2\) from the geometric mean, and \(h\) from the
amplitude. Their sum is \(5h/2\le3h\). After substituting
\(H_Z\le P\sigma_Z^4\), either orientation leaves exponent at most
\[
   12h+C_2m^2/2+|C_2|m^4-a\le-12p^2
   \qquad(Z=X,Y).
\]
There are two choices of the side giving the gain, and
\(\sum_{X,Y}\phi(X)\phi(Y)\le\Phi^2\) by
Lemma~\ref{lem:gram-label-sum}.
\end{proof}

\paragraph{Choosing a common moment order.}

The three counts impose different lower bounds on the order of the moment. The local estimates also impose an upper bound, so one must exhibit a single admissible integer for all three.

\begin{lemma}[A common moment order for the three counts]\label{lem:row-moment-budget}
Put \(p=m/2=r+1\) and \(h=cr^2\). Suppose the size parameters satisfy
\[
 m+3\le k,\qquad 2+\log_2 f\le m,\qquad 3|C_2|m^5\le k .
\]
Let \(a_0\) be the largest of the two lower bounds in
Proposition~\ref{prop:row-higher-corank} and the lower bound in
Proposition~\ref{prop:summed-row-corank-one}. All three counts admit the same integral moment order:
\[
 q=\lceil p a_0\rceil\le(9c+3)km^2,\qquad a=q/p\ge a_0 .
\]
\end{lemma}

\begin{proof}
Let \(a_0\) be the maximum of the three required bounds,
and put \(q=\lceil p a_0\rceil\). With \(p=m/2\) and \(h=cr^2\),
use \(m+3\le k\) and \(2+\log_2 f\le m\). In the higher
corank bounds, \(p(4ck+4E_k)\le 5ckm^2\) and
\(4ph\le ckm^2/2\). In the corank-one bound,
\(12ph\le 3ckm^2/2\).

In all cases
\[
 p|C_2|(m^4+m^2/2)\le k/3,\qquad 12p^3\le 3km^2/2
\]
under \(3|C_2|m^5\le k\). Adding one for the ceiling leaves
\(q\le(9c+3)km^2\). The last condition on \(k\) holds
eventually by the selected size bounds.

In particular the actual endpoint radius, for every minor order
\(l\le p=m/2\), is bounded by
\[
  ql\le pq\le \tfrac12(9c+3)km^3 .
\]
Fix \(C_n=(9c+3)/2+1\) in Lemma~\ref{lem:minor-orbital-shift}.
At depth \(j<q\) the one-step proof also needs the moment at
\(\nu+\mathbf1_I\), whose radius is \((j+1)l\le pq\).
The internal moment computation needs
\[
      p(24r^3+1)+pq\le km^5,\qquad p(24r^3+1)\le2n^2 .
\]
Both follow eventually from fixed \(c\), \(p=m/2\), and \(k\ge m^4\).
The marked gap dominates the additional \(200pq\) needed for
point evaluation since \(\gamma=nm^{10}\), \(n\ge k\).
Thus all endpoints lie in the same chamber and the same quantitative
window for both words.

In each case the two label sums contribute one factor \(\Phi^2\)
by Lemma~\ref{lem:gram-label-sum}. Together with
\(P^{2(E_s+h)}\) at higher corank, or \(P^{6h}\) at corank one,
this contributes
\(b^{O_{c,C_h}((s+m)m^4f)}\) at higher corank
and \(b^{O_{c,C_h}(m^5f)}\) at corank one. The fixed
constants in these powers are uniform in the labels,
the places, and the group size. At higher corank the additional
factor \(\Xi_s(P^2)\) remains in \eqref{eq:high-corank-summed-prices}.
\end{proof}

\subsection{The generic corank-one sum on an equal-word block}

For two copies of a word the allowed error is larger than in the
comparison of different words. The generic partner count in the two
orders already suffices for this diagonal estimate. It includes
pairs satisfying Definition~\ref{def:corank-one-scalar}.

\begin{auxlemma}[Generic corank-one pairs on a diagonal block]\label{lem:equal-word-generic-corank-one}
Fix a marked word \(u\) and its prime \(T=T_u\). Take both words
equal to \(u\) in the amplitude and height estimates of this
section, with \(q=\lceil p a_0\rceil\) as in
Lemma~\ref{lem:row-moment-budget}. Let \(\mathcal S\) be any subset
of the matching ordered pairs of generic terms having positive
single-term Gram invariants and common rank \(k-1\). Write
\(w_{t,t'}\) for the counting weight of Lemma~\ref{lem:row-fibre-weight},
\(a_t\) for the archimedean coefficient norm there, and \(\Gamma_1\)
for the finite diagonal squared norm of Lemma~\ref{lem:higher-corank-total-error}.
For \(L=\log b\), \(M=m^5fL\), there is a constant \(C_B\)
independent of \(C_\Gamma\) such that, on all sufficiently late rows,
\begin{equation}\label{eq:equal-word-generic-corank-one}
 \Gamma_1\sum_{(t,t')\in\mathcal S} w_{t,t'}|a_ta_{t'}|
             \le \exp(C_B M)\,\theta_{uu},
 \qquad \theta_{uu}=d_u^2/D_* .
\end{equation}
\end{auxlemma}
\begin{proof}
\proofstep{Counting with either term first.}
Fix reduced labels \(X,Y\) for the two terms. In the count of
Proposition~\ref{prop:flag-pair-counts}, the orthogonal determinant
of a generic first term with label \(X\) is \(G_X\). By
\eqref{eq:generic-first-term-factor} its normalized genus mean is at most
\[
 H_X^h\left(\frac{e\pi}{k\upsilon}\right)^{kf}
          (G_X/N_T)^k=H_X^h n_X .
\]
The common rank is \(k-1\), so the integral orthogonal lattice of
the first term lies in its partner modulo the isotropic space.
After fixing that lattice, \eqref{eq:flag-over-P} bounds the number
of partner classes with label \(Y\) by \(c^f(H_XH_Y)^h\), with
\(h=cr^2\) for the common counting constant fixed before \(q\).
The level congruence makes the comparison isometry unique in the
pair class, as in Proposition~\ref{prop:flag-pair-counts}.

Switching the two terms replaces the comparison isometry by its
inverse. Both terms are generic, and the base row retains the same
flag, prime \(T_u\) and reciprocal stabilizer weight in this
count. For the total row mass \(M_u=m_k(T_u^k-1)^f\) it follows that
\[
 \frac1{M_u}\sum_{\substack{(t,t')\in\mathcal S\\(t,t')\mapsto(X,Y)}}w_{t,t'}
 \le c^f(H_XH_Y)^h\min(H_X^h n_X,H_Y^h n_Y).
\]
These counts depend only on integrality and on the common rank. In
particular the same estimate holds on restricting \(\mathcal S\)
to the scalar pairs of Definition~\ref{def:corank-one-scalar}.

\proofstep{Keeping the height and the column dimension separate.}
Set \(D_L=C+6CC_h\). In the proof of Lemma~\ref{lem:row-height-products}
the first product is bounded by
\(b^{Cm^4f}(B^f\sigma_Z)^{6Cm^2}\).
Using this expression in the derivation of \eqref{eq:row-l2}
gives the squared amplitude bound for a term of label \(Z\),
\[
 n_Z|a_t|^2/d
 \le H_Z^{2h}\,b^{D_Lm^4f}\,\tau\,
                   \sigma_Z^{C_2m^2-2a},\qquad
 \tau=(1+\delta)d_u/D_*,\quad d=\dim\pi_u .
\]
The finite levels enter \(G_Z\) and \(n_Z\) in the same way in
both calculations. Dividing the minimum in the fibre count by
\(\sqrt{n_Xn_Y}\) costs at most \((H_XH_Y)^{h/2}\), by the
geometric mean. The two square roots of the amplitude bounds then
leave the height power \(5h/2\) on each label.

Now \(H_Z\le P\sigma_Z^4\), and our moment order satisfies
\[
 10h+C_2m^2/2-a\le-12p^2 .
\]
Thus both remaining label weights sum by Lemma~\ref{lem:gram-label-sum}.
Use the finite norm identity \(\Gamma_1 M_u d=d_u\) from
Lemma~\ref{lem:higher-corank-total-error}. The sum on the left of
\eqref{eq:equal-word-generic-corank-one} is at most
\[
       (1+\delta)c^f P^{5h}\Phi^2 b^{D_Lm^4f}\,\theta_{uu}.
\]
This is an absolute majorant, so it also bounds the sum of the real
parts of the matching coefficient inner products.

\proofstep{Size and order of the constants.}
Write \(A_P=c^2+2\max(C_h-c,0)\) and take \(C_*\ge1\) in the
packing bound. Since \(\Phi=(C_*^fP^3)^{p^2}\), the logarithm
of the multiplier is
\[
 \log(1+\delta)+f\log c+2p^2f\log C_*
 +(5h+6p^2)A_Pm^3fL+D_Lm^4fL .
\]
For \(L\ge1,\ \delta\le1\), use \(p\le m\) and \(h\le cm^2\).
One may take
\[
 C_B=1+\log c+2\log C_*+(5c+6)A_P+D_L .
\]
The counting exponent \(c\), the local bounds \(C,C_h,C_2\), and
the moment constants in Lemma~\ref{lem:row-moment-budget} are
fixed by the level, weight and measure bounds. In particular the
upper bound \(\log T_u\ll mL\) already follows from
\(T_u\le(1+1/(fk))X_u\) in Proposition~\ref{prop:prime-target-size}.

After \(C_\Gamma\) is chosen, the actual target has the correction
\(2C_\Gamma m^5L/k\) in its logarithm. This tends to zero since
\(k\ge m^5Lx^2\). The inequality \(11f2^{-k}\le\log(1+b^{-7})\)
also holds eventually. They supply the selected prime range and
\(\delta=b^{-7}\) at a possibly later threshold. The value of
\(C_B\) above applies at that threshold.
\end{proof}
\subsection{The diagonal pairs with a nonprimitive frame span}

The partner count used for a generic term fixes its full integral
orthogonal lattice. A nonprimitive pair of frames can have different
integral orthogonal lattices even when the rational orthogonal spaces
agree modulo the flag. Retaining the two indices in the common
congruence lattice gives the required bound at corank one.

\begin{auxlemma}[Corank one with a nonprimitive frame span]\label{lem:nonprimitive-diagonal}
Fix a marked word \(u\) on the row and use the moment order of
Lemma~\ref{lem:row-moment-budget}. In the matching ordered pairs on
the block \((u,u)\), let \(\mathcal S\) be any subset whose two
single-term Gram invariants are positive and whose common rank is
\(k-1\). With the weights, finite norm and coefficient norms of
Lemma~\ref{lem:equal-word-generic-corank-one},
\begin{equation}\label{eq:nonprimitive-diagonal}
  \Gamma_1\sum_{(t,t')\in\mathcal S}w_{t,t'}|a_t a_{t'}|
       \le \exp(C_N m^5f\log b)\,\theta_{uu}.
\end{equation}
Here \(C_N\) depends only on the fixed counting, level and local
amplitude constants and can be chosen before \(C_\Gamma\).
In particular the bound applies to the pairs in the nongeneric
term group on the diagonal.
\end{auxlemma}

\begin{proof}
\proofstep{The indices of one marked pair of frames.}
Write \(R=\mathcal O_E\), \(g=g_0g_1\), \(p=r+1\), and
\(N=T_u^f\). In an integral unimodular row lattice \(\Lambda\)
denote the fixed flag by \((X,z)\) and its two partners by
\((X_j,z_j)\), \(j=1,2\). Here \(X,X_j\) are the isotropic spaces
and \(z,z_j\) the distinguished vectors. Put
\[
 \begin{aligned}
 S_j&=(\Lambda\cap X)+Rz+(\Lambda\cap X_j)+Rz_j,&
 W_j&=E S_j,\\
 B_j&=(\Lambda\cap W_j)/S_j,&
 U_j&=\Lambda\cap W_j^\perp,\qquad
 \beta_j=|B_j|/g^{mf}.
 \end{aligned}
\]
The two augmented flag spaces in \(W_j\) form a direct sum by
positivity of the single-term Gram invariant. Let \(G_j\) be the
normalized determinant \(G_Z\) of the corresponding reduced label.
The determinant calculation in Proposition~\ref{prop:flag-pair-counts}
before taking the primitive hull gives \(d(S_j)=g^{mf}G_j\).
The saturated lattice \(\Lambda\cap W_j\) and \(U_j\) have
opposite discriminant pairings of equal order by unimodularity.
Thus in the determinant convention \(d(L)=\sqrt{|L^\#/L|}\),
\[
   d(U_j)=d(\Lambda\cap W_j)
          =\frac{d(S_j)}{|B_j|}=\frac{G_j}{\beta_j}.
\]

At primes over \(g\), the marked flags have unit integral bases and
their level comparison is upper unitriangular. Dividing the \(p\)
corresponding frame differences by \(g\) gives integral classes
in \(B_j[g]\). If a combination is in \(S_j\), projection onto the
first augmented flag in the rational direct sum puts its coefficients,
after the unitriangular change of basis, in \(gR\) locally.
The distinguished negative vector is primitive in the middle
quotient \((\Lambda\cap X^\perp)/(\Lambda\cap X)\), so its
coordinate obeys the same condition.
Conversely multiples of \(g\) give zero classes. Hence the subgroup
generated locally has the size of \((R/gR)^p\) globally. Equivalently,
adjoining these marked divisions to \(S_j\) gives an intermediate
lattice of index \(g^{2fp}=g^{mf}\). In particular \(\beta_j\) is
the residual absolute saturation index and
\[
       |gB_j|=|B_j|/|B_j[g]|\le\beta_j .
\]

For \(V=X\) or \(X_j\), write \(L(V)=\Lambda\cap V^\perp\) and
\(V_R=\Lambda\cap V\). Define the enlarged orthogonal lattice
\(A_j\subset W_j^\perp\) by the existence of representatives
\[
 \ell\in L(X),\qquad \ell_j\in L(X_j),\qquad
 x-\ell\in X,\quad x-\ell_j\in X_j,\quad
 \ell-\ell_j\in g\Lambda+X_R+(X_j)_R .
\]
These conditions are \(R\)-linear. The map
\(x\mapsto[\ell-\ell_j]\in B_j\) is independent of representatives,
since two choices on each side differ by an integral flag vector.
Its image belongs to \(gB_j\): subtract the two integral flag
vectors in the congruence and divide by \(g\). The resulting vector
lies in \(\Lambda\cap W_j\).

To identify the kernel, write \(\ell-\ell_j=v+v_j\) in \(S_j\),
with \(v,v_j\) integral in their respective augmented flag spans.
Then
\[
       x-\ell+v=x-\ell_j-v_j
        \ \in\ (X+Ez)\cap(X_j+Ez_j)=0 .
\]
It follows that \(x=\ell-v\in U_j\). Conversely an integral
\(x\in W_j^\perp\) permits the representatives \(\ell=\ell_j=x\).
We have proved
\begin{equation}\label{eq:nonprimitive-one-index}
 I_j=[A_j:U_j]\le\beta_j,\qquad
 d(A_j)I_j^2=d(U_j)I_j\le G_j .
\end{equation}

\proofstep{The integral lattice shared by a matching pair.}
Write \(\eta\) for its integral comparison isometry, carrying the
second flag of the first term to that of the second. The joint
rank is the dimension of the correspondence
\[
 x\in W_1^\perp,\quad y\in W_2^\perp,\qquad
 x-y\in X,\quad \eta x-y\in X_2 .
\]
Each projection is injective since \(W_j^\perp\cap X=0\).
Both target spaces have dimension \(k-1\), so the correspondence
gives a unique \(y\) for each \(x\), and conversely. Projection
modulo \(X\) identifies these rational spaces isometrically.

The marking at the common finite index requires, for every
\(\ell\in L(X)\) and \(\ell_1\in L(X_1)\),
\[
 \ell-\ell_1\in g\Lambda+X_R+(X_1)_R
 \quad\Longleftrightarrow\quad
 \ell-\eta\ell_1\in g\Lambda+X_R+(X_2)_R .
\]
For \(x\in A_1\), keep its first representative \(\ell\) and
use \(\eta\ell_1\) as the second representative of \(y\).
Indeed \(y-\ell\in X\) and
\(y-\eta\ell_1=(y-\eta x)+\eta(x-\ell_1)\in X_2\).
The displayed congruence gives \(y\in A_2\).
The reverse implication and \(\eta^{-1}\) give the converse.
Thus \(A_1,A_2\) have the same projection \(A\) in the
positive definite subspace \(z^\perp\) of \(X^\perp/X\).

In that quotient let \(P_0\) be the intersection of the
projections of \(U_1,U_2\). These are integral sublattices of
the definite middle lattice of the row. The residue map into
\(A/U_1\oplus A/U_2\), with the lattices understood by projection,
now gives
\[
 \begin{aligned}
  {}[A:P_0]&\le I_1I_2,\\
  d(P_0)&\le d(A)I_1I_2
      =\sqrt{d(A_1)I_1^2\,d(A_2)I_2^2}
      \le \sqrt{G_1G_2}.
 \end{aligned}
\]
This uses the level equivalence in both directions. It does not
require either individual frame span to be primitive.

\proofstep{Counting the common sublattices in the row measure.}
Fix reduced labels \(Z_1,Z_2\), write \(H_j=H_{Z_j}\),
\(H=H_1H_2\), \(D=\sqrt{G_1G_2}\), and set \(\lambda=\log_2(2f)\).
The selected prime satisfies \(N=T_u^f\le b^{C_T mf}\le H\),
and Lemma~\ref{lem:reduced-flag-heights} gives \(D\le H\). We may take \(D\ge1\)
whenever there is a pair. For a fixed \(P_0\) with \(d(P_0)\le D\),
Proposition~\ref{prop:integral-flag-frames} applies separately to
the two frames with \(s=1\) and \(Q_j=H_j^2\).
Their classes determine the marked pair class by the uniqueness
argument of Proposition~\ref{prop:flag-pair-counts} (the difference
is a finite-order isometry congruent to the identity modulo \(g\)).
The two bounds in \eqref{eq:flag-over-P} therefore cost together
at most
\[
             9^{2f}H^{8r^2+18r+12}.
\]

If \(P_0\) saturates to determinant \(d_0\), then
\([\overline P_0:P_0]\le D/d_0\). The index count in
Proposition~\ref{prop:primitive-flag-means} bounds all possibilities
over this saturation by \((D/d_0)^{k+\lambda+1}\).
The corank-one part of the same proposition gives the mean count
of saturations under the reciprocal row stabilizer measure.
The two large powers combine as
\[
 (d_0/N)^{k-1}(D/d_0)^{k+\lambda+1}
    =(D/N)^k(N/D)(D/d_0)^{\lambda+2}.
\]

There are at most \(1+D^2\) determinant values since \(d_0^2\)
is an integer. The remaining factors in this average are bounded by
\[
  h_E(4k\sqrt{\mathrm{rd}(E)})^f
        H^{11+3\log_2 f}(1+D^2)(N/D)D^{\lambda+2}.
\]
On the marked range, \((kb)^8\le g\le b^{C_g}\),
\(\mathrm{rd}(E)\le b\), and
\(h_E\le\mathrm{rd}(E)^{f(1+\lambda)}\).
Since \(H_j\ge b^{c^2m^3 f}\) and \(2+\log_2 f\le m\), these
factors and the two frame counts cost at most \(H^{e_0}\) for
\(e_0=C_0m^2\) with \(C_0\) fixed independently of \(C_\Gamma\).

Restoring the remaining corank-one mean factor yields
\[
 \frac{1}{M_u}\sum_{\substack{(t,t')\in\mathcal S\\
                     (t,t')\mapsto(Z_1,Z_2)}}w_{t,t'}
 \le H^{e_0}\left(\frac{e\pi}{k\upsilon}\right)^{kf}
                  (D/N)^k
     = H^{e_0}\sqrt{n_{Z_1}n_{Z_2}},
\]
where \(M_u=m_k(T_u^k-1)^f\)
is the total mass from Proposition~\ref{prop:flag-row-mass}.
The same weighted mean includes every finite coarse row class
of the specified pair subset.

\proofstep{Restoring the archimedean coefficient and the finite norm.}
Use the squared-norm bound with \(D_L=C+6CC_h\) from the proof of
Lemma~\ref{lem:equal-word-generic-corank-one}. Its proof applies to
every positive single-term Gram invariant. For \(d=\dim\pi_u\) it reads
\[
 |a_t|/\sqrt d
 \le H_j^h\,b^{D_Lm^4 f/2}\sqrt{\tau/n_{Z_j}}\,
                         \sigma_{Z_j}^{C_2m^2/2-a}.
\]

Thus the square roots of the two \(n_{Z_j}\) cancel on this very
fibre. After \(H_j\le P\sigma_{Z_j}^4\), the exponent of each
\(\sigma_{Z_j}\) is \(4(e_0+h)+C_2m^2/2-a\).
The actual order \(a=\lceil pa_0\rceil/p\) of
Lemma~\ref{lem:row-moment-budget} satisfies \(a\ge4ck\).
Since \(k\ge m^5(\log b)x^2\), eventually
\[
        a\ge4(e_0+h)+|C_2|m^2/2+12p^2 .
\]

Lemma~\ref{lem:gram-label-sum} now bounds the two independent
label sums by \(\Phi^2\). With
\(\Gamma_1M_ud=d_u\) and \(\tau=(1+\delta)d_u/D_*\),
the resulting bound is
\[
     (1+\delta)P^{2(e_0+h)}\Phi^2 b^{D_Lm^4f}\,\frac{d_u^2}{D_*}.
\]
Indeed the weights just counted are the finite row weights, while
the two coefficient norms use the real orbital Haar of
Proposition~\ref{prop:two-row-amplitude-bounds}.
Their only remaining finite multiplier is \(\Gamma_1\).
Here \(\log P=O(m^3f\log b)\), \(e_0+h=O(m^2)\), and
\(\Phi=(C_*^fP^3)^{p^2}\), proving \eqref{eq:nonprimitive-diagonal}.
\end{proof}

